
\subsection{Ambient positive curvature near the same-density teacher}
\label{subsec:post-escape-convexity-on-the-same-density}

This subsection proves positive spherical product pushforward second variation
at, and near, the exact-fit same-density teacher.  It does not assert that an
escaping gradient-flow trajectory reaches this neighborhood, or that the
teacher is stationary for the regularized objective.

The descent direction in \cref{lemma:descent-direction-antipodal-dirac-from-unweighted-imbalance} creates a same-sign odd mode. This motivates the same-density branch
\begin{equation}
    \mathcal M_+
    \coloneq
    \left\{
        \mu_1\otimes \mu_2
        \mid
        \mu_1=\mu_2
    \right\}.
\end{equation}
On $\mathcal M_+$ we have $A_n=0$ for every $n\in \mathcal O$ and $S_n=T_{1,n}=T_{2,n}$. By \cref{prop:reference-fit-minimum-on-Mplus}, the reference zero-fit point of the data term lies on this branch.

\subsubsection{Invariance of the same-density branch}

We first check that the same-density branch is dynamically consistent with the product WGF.
\begin{proposition}[Invariance of $\mathcal M_+$ under the antipodal metric objective]
\label{lem:invariance-same-density-branch}
\textbf{Assumptions used.} \cref{assumption:on-sphere,assumption:odd-symmetric-teacher}, together with the exchange symmetry of $F$ in $(\mu_1,\mu_2)$.

If \(\mu_0=\mu_{1,0}\otimes \mu_{2,0}\in \mathcal M_+\) and \(\mu_t\) is the
unique corresponding product WGF solution on an interval \([0,T)\) in the
fixed-assignment cap class \(\mu_{i,t}\in\mathcal P_\pm(r_{\mathrm{cap}})\), \(i=1,2\), then
\begin{equation}
    \mu_t\in \mathcal M_+
    \qquad
    \forall t\in[0,T).
\end{equation}
\end{proposition}

\begin{proof}
\noindent\textbf{Step 1: exchange the two factors in the objective.}
If $\mu_1=\mu_2$, then $T_{1,n}=T_{2,n}$ for every $n\in \mathcal O$. 

The fitting part is invariant under interchanging the two factors because $\mu_1^o=\mu_2^o$. 

The metric regularizer is exchange-invariant because it depends on the two factors only through
\[
    Z_\pm(\mu_1,\mu_2)
    =
    W_2^2(\mu_1,\nu_\star^\pm)+W_2^2(\mu_2,\nu_\star^\pm).
\]

\par\noindent\textbf{Step 2: transform the characteristics and apply uniqueness.}
In particular, exchange symmetry gives
\[
 F(\mu_1,\mu_2)=F(\mu_2,\mu_1),\qquad
 g_1(\mu_1,\mu_2;w)=g_2(\mu_2,\mu_1;w),
\]
where \(g_i=\nabla_\Sphere(\delta F/\delta\mu_i)\), as in
\cref{eq:mf-branch-radial-notation}. For the swapped pair
\(\widetilde\mu_{1,t}=\mu_{2,t}\), \(\widetilde\mu_{2,t}=\mu_{1,t}\),
its first continuity equation is
\[
 \partial_t\widetilde\mu_{1,t}
 -\nabla_\Sphere\cdot
  (g_1(\widetilde\mu_{1,t},\widetilde\mu_{2,t};\cdot)\widetilde\mu_{1,t})
 =\partial_t\mu_{2,t}
 -\nabla_\Sphere\cdot(g_2(\mu_{1,t},\mu_{2,t};\cdot)\mu_{2,t})=0.
\]
The other equation follows by exchanging indices. Since
\(\mu_{1,0}=\mu_{2,0}\), the swapped pair has the same initial state.
Uniqueness on the stated interval gives
\(\widetilde\mu_{i,t}=\mu_{i,t}\), hence \(\mu_{1,t}=\mu_{2,t}\).
\end{proof}

\subsubsection{Positive curvature at and near the same-density teacher}

Let
\begin{equation}
 \nu^o:=\mu_1^o=\mu_2^o,
 \qquad
 \mu_+^o:=\nu^o\otimes\nu^o\in\mathcal M_+.
\end{equation}
At this state the student output equals the teacher output, so the fitting
residual vanishes. The next proposition combines the resulting nonnegative
fitting curvature with the strict metric curvature and shows that the margin
persists locally.

\begin{proposition}[Ambient positive curvature near a same-density teacher]
\label{prop:Mplus-teacher-local-positive-curvature}
\textbf{Assumptions used.}
\cref{assumption:on-sphere,assumption:odd-symmetric-teacher}.
Suppose
\[
 \nu^o\in\mathcal P_\pm(r_{\mathrm{cap}}),
 \qquad
 Z_o:=Z_\pm(\mu_+^o)
 =2W_{2,\Sphere}^2(\nu^o,\nu_\star^\pm)>0.
\]
Then there exist \(\rho_o>0\) and \(c_o>0\), depending only on the fixed
model parameters and the teacher, such that the following holds. For every
\(\mu=\mu_1\otimes\mu_2\in\mathcal X_\pm(r_{\mathrm{cap}})\) satisfying
\[
 W_{2,\Sphere}^2(\mu_1,\nu^o)
 +W_{2,\Sphere}^2(\mu_2,\nu^o)<\rho_o^2,
\]
and every smooth product tangent field \(v=(v_1,v_2)\) whose perturbation
paths preserve the fixed-assignment caps for small parameter values,
\begin{equation}
\label{eq:mf-teacher-local-positive-curvature}
 D_{W_2}^2F(\mu)[v,v]
 \ge c_o\|v\|_{L^2(\mu;\mathbb R^{d_w})}^2.
\end{equation}
No density assumption on \(\nu^o\) is required.
\end{proposition}

\begin{proof}
\noindent\textbf{Step 1: positive curvature at the exact-fit teacher.}
For smooth tangent fields \(v_i\) on the two factors, let
\[
 R_{i,h}(w):=\frac{w+hv_i(w)}{\|w+hv_i(w)\|_{\mathbb R^{d_w}}},
 \qquad
 \nu_{i,h}:=(R_{i,h})_\#\nu^o,
 \qquad
 \mu_h:=\nu_{1,h}\otimes\nu_{2,h}.
\]
The normalized retraction has the same second jet as the spherical EXP path
by \cref{eq:mf-retraction-second-jet}, so its second derivative computes
\(D_{W_2}^2F(\mu_+^o)[v,v]\). Set
\[
 Q(h,z):=h_{\nu_{1,h}}(z)h_{\nu_{2,h}}(z)^\top,
 \qquad
 B(z):=h_{\nu^o}(z)h_{\nu^o}(z)^\top.
\]
Exact fitting means \(Q(0,z)=B(z)\). Therefore, differentiating the squared
residual twice gives
\begin{align}
 D_{W_2}^2F_{\mathrm{fit}}(\mu_+^o)[v,v]
 &=\left.\frac{\dd^2}{\dd h^2}\right|_{h=0}
   \frac12\|Q(h)-B\|_{L^2_z(\mathrm F)}^2\notag\\
 &=\|\dot Q(0)\|_{L^2_z(\mathrm F)}^2\notag\\
 &=\sum_{p\ge0}\left\|
   \left.\frac{\dd}{\dd h}\right|_{h=0}
   \widetilde{\mathcal M}_p(\Theta(\mu_h))
   \right\|_{\mathrm{HS}}^2\ge0.
 \label{eq:mf-teacher-fitting-square}
\end{align}
The last equality is the finite active Hermite expansion from
\cref{lemma:EMN-expansion-F,eq:ambient-active-fit-functional}.

By \cref{lem:W2-power-regularizer-Hessian-on-fixed-caps}, the metric term at
the same state satisfies
\[
 D_{W_2}^2F_{\mathrm{reg}}(\mu_+^o)[v,v]
 \ge\lambdaReg\kappa_{\mathrm{cap}}Z_o^{(p_W-2)/2}
      \|v\|_{L^2(\mu_+^o;\mathbb R^{d_w})}^2.
\]
The coefficient is positive because \(Z_o>0\). Adding this inequality to
\cref{eq:mf-teacher-fitting-square} proves strict ambient curvature at the
teacher itself.

\par\noindent\textbf{Step 2: a lower bound at a nearby cap-supported state.}
Fix \(\mu=\mu_1\otimes\mu_2\in\mathcal X_\pm(r_{\mathrm{cap}})\), and use the
same normalized retractions based at each \(\mu_i\). Write
\[
 h_i(h,z):=\int\sigma(R_{i,h}(w)^\top z)\,\dd\mu_i(w),
 \qquad
 V_i:=\|v_i\|_{L^2(\mu_i;\mathbb R^{d_w})}.
\]
Dots below denote derivatives at \(h=0\). From
\cref{eq:mf-retraction-second-jet},
\(\dot R_{i,0}=v_i\) and
\(\ddot R_{i,0}=-\|v_i\|_{\mathbb R^{d_w}}^2w\), so differentiation under
the integral gives
\begin{equation}
\label{eq:mf-teacher-output-derivatives}
\begin{aligned}
 \dot h_i(z)
 &=\int\sigma'(w^\top z)(v_i(w)^\top z)\,\dd\mu_i(w),\\
 \ddot h_i(z)
 &=\int\left[
   \sigma''(w^\top z)(v_i(w)^\top z)^2
   -\|v_i(w)\|_{\mathbb R^{d_w}}^2
      \sigma'(w^\top z)(w^\top z)
   \right]\dd\mu_i(w).
\end{aligned}
\end{equation}
When \(v_i(w)\ne0\), the vectors
\(w\) and \(v_i(w)/\|v_i(w)\|_{\mathbb R^{d_w}}\) are orthogonal unit
vectors because \(v_i(w)\in T_w\Sphere^{d_w-1}\). Hence their projections
of \(z\sim N(0,I_{d_w})\) are independent standard Gaussians. Substitution
into the definitions of \(M_4,K_1,K_2\) in
\cref{eq:mf-gaussian-derivative-constants}, followed by Minkowski and
Cauchy--Schwarz, yields
\begin{equation}
\label{eq:mf-teacher-output-derivative-bounds}
 \|h_i(0)\|_{L^4_z}\le M_4,
 \qquad
 \|\dot h_i\|_{L^4_z}\le K_1V_i,
 \qquad
 \|\ddot h_i\|_{L^4_z}\le K_2V_i^2.
\end{equation}
At \(v_i(w)=0\), the corresponding integrands in
\cref{eq:mf-teacher-output-derivatives} vanish, so the same bounds hold
without division by zero. The finite polynomial activation and compact
sphere justify the differentiations and Gaussian integration.

For
\(Q(h,z):=h_1(h,z)h_2(h,z)^\top\), the product rule gives
\[
 \dot Q=\dot h_1h_2^\top+h_1\dot h_2^\top,
 \qquad
 \ddot Q=\ddot h_1h_2^\top+2\dot h_1\dot h_2^\top+h_1\ddot h_2^\top.
\]
Since \(\|xy^\top\|_{\mathrm F}=\|x\|_2\|y\|_2\), H\"older and
\cref{eq:mf-teacher-output-derivative-bounds} give, term by term,
\[
\begin{aligned}
 \|\ddot h_1h_2^\top\|_{L^2_z(\mathrm F)}&\le M_4K_2V_1^2,\\
 \|2\dot h_1\dot h_2^\top\|_{L^2_z(\mathrm F)}&\le2K_1^2V_1V_2,\\
 \|h_1\ddot h_2^\top\|_{L^2_z(\mathrm F)}&\le M_4K_2V_2^2.
\end{aligned}
\]
Using \(2V_1V_2\le V_1^2+V_2^2\), we obtain
\begin{equation}
\label{eq:mf-teacher-second-output-bound}
 \|\ddot Q\|_{L^2_z(\mathrm F)}
 \le (M_4K_2+K_1^2)
       \|v\|_{L^2(\mu;\mathbb R^{d_w})}^2.
\end{equation}

Recall \(B(z)=h_{\nu^o}(z)h_{\nu^o}(z)^\top\), and set
\[
 E(\mu):=\|h_{\mu_1}h_{\mu_2}^\top-B\|_{L^2_z(\mathrm F)},
 \qquad C:=M_4K_2+K_1^2.
\]
The constant \(C\) is proof-local and depends only on the fixed activation.
Differentiating
\(F_{\mathrm{fit}}(\mu_h)=\frac12\|Q(h)-B\|_{L^2_z(\mathrm F)}^2\),
dropping the nonnegative \(\|\dot Q\|_{L^2_z(\mathrm F)}^2\) term, and using
\cref{eq:mf-teacher-second-output-bound} gives
\begin{align*}
 D_{W_2}^2F_{\mathrm{fit}}(\mu)[v,v]
 &=\|\dot Q\|_{L^2_z(\mathrm F)}^2
   +\langle Q(0)-B,\ddot Q\rangle_{L^2_z(\mathrm F)}\\
 &\ge-CE(\mu)\|v\|_{L^2(\mu;\mathbb R^{d_w})}^2.
\end{align*}
Adding the regularizer estimate in
\cref{eq:W2-power-Hessian-lower-bound-fixed-caps} yields the key nearby
bound
\begin{equation}
\label{eq:mf-teacher-nearby-curvature-bound}
 D_{W_2}^2F(\mu)[v,v]
 \ge\left[
   \lambdaReg\kappa_{\mathrm{cap}}Z_\pm(\mu)^{(p_W-2)/2}
   -CE(\mu)
 \right]\|v\|_{L^2(\mu;\mathbb R^{d_w})}^2.
\end{equation}

\par\noindent\textbf{Step 3: keep the coefficient positive in a teacher neighborhood.}
Let
\[
 d:=\left[
 W_{2,\Sphere}^2(\mu_1,\nu^o)+
 W_{2,\Sphere}^2(\mu_2,\nu^o)
 \right]^{1/2}.
\]
For sphere points \(w,y\), let \(c:[0,1]\to\Sphere^{d_w-1}\) be a shortest
constant-speed arc from \(y\) to \(w\). Then
\(\|c'(s)\|_{\mathbb R^{d_w}}=d_\Sphere(w,y)\) and
\(c(s)^\top c'(s)=0\). Thus, when the speed is nonzero,
\(c(s)\) and \(c'(s)/\|c'(s)\|_{\mathbb R^{d_w}}\) are orthogonal unit
vectors. The same Gaussian calculation used in
\cref{eq:mf-teacher-output-derivative-bounds} gives
\[
 \|\sigma'(c(s)^\top z)(c'(s)^\top z)\|_{L^4_z}
 =K_1\|c'(s)\|_{\mathbb R^{d_w}}.
\]
The identity is immediate when the speed is zero. Integrating along the arc
and applying Minkowski gives
\[
 \|\sigma(w^\top z)-\sigma(y^\top z)\|_{L^4_z}
 \le K_1d_\Sphere(w,y).
\]
For an optimal coupling \(\pi_i\) of \(\mu_i\) and \(\nu^o\), another use
of Minkowski and Cauchy--Schwarz therefore yields
\[
\begin{aligned}
 \|h_{\mu_i}-h_{\nu^o}\|_{L^4_z}
 &\le K_1\int d_\Sphere(w,y)\,\dd\pi_i(w,y)\\
 &\le K_1W_{2,\Sphere}(\mu_i,\nu^o).
\end{aligned}
\]
Subtracting the two product outputs,
\[
 h_{\mu_1}h_{\mu_2}^\top-B
 =(h_{\mu_1}-h_{\nu^o})h_{\mu_2}^\top
 +h_{\nu^o}(h_{\mu_2}-h_{\nu^o})^\top.
\]
Since \(\|h_{\mu_i}\|_{L^4_z},\|h_{\nu^o}\|_{L^4_z}\le M_4\) by
\cref{eq:mf-output-fourth-moment-bound}, H\"older and the preceding feature
bound give \(E(\mu)\le\sqrt2M_4K_1d\). Moreover,
\[
 d=W_{2,\mathcal M}(\mu,\mu_+^o),\qquad
 \sqrt{Z_\pm(\mu)}=W_{2,\mathcal M}(\mu,\mu^\dagger),\qquad
 \sqrt{Z_o}=W_{2,\mathcal M}(\mu_+^o,\mu^\dagger),
\]
by the product metric in \cref{eq:product-Wasserstein-metric} and the
definition of $Z_\pm$ in \cref{eq:fixed-assignment-W2-square-formula},
together with the definition of $Z_o$ in the proposition statement.
Therefore the reverse triangle inequality gives
\(\left|\sqrt{Z_\pm(\mu)}-\sqrt{Z_o}\right|\le d\). Combining the two
estimates yields
\begin{equation}
\label{eq:mf-teacher-neighborhood-continuity}
 E(\mu)\le\sqrt2M_4K_1d,
 \qquad
 \left|\sqrt{Z_\pm(\mu)}-\sqrt{Z_o}\right|\le d.
\end{equation}

As \(d\to0\), \cref{eq:mf-teacher-neighborhood-continuity} shows that
\(E(\mu)\to0\) and \(Z_\pm(\mu)\to Z_o\). Hence the bracket in
\cref{eq:mf-teacher-nearby-curvature-bound} converges to the positive number
\(\lambdaReg\kappa_{\mathrm{cap}}Z_o^{(p_W-2)/2}\). Choose
\(\rho_o>0\) so that, whenever \(d<\rho_o\), the regularizer coefficient is
at least three quarters of this limit and \(CE(\mu)\) is at most one quarter
of it. Then \cref{eq:mf-teacher-nearby-curvature-bound} gives
\cref{eq:mf-teacher-local-positive-curvature} with
\[
 c_o:=\frac12\lambdaReg\kappa_{\mathrm{cap}}Z_o^{(p_W-2)/2}>0.
\]
The neighborhood is relative to \(\mathcal X_\pm(r_{\mathrm{cap}})\): small
Wasserstein distance alone does not force the entire support to remain in the
fixed-assignment caps. None of the estimates uses a teacher density.
\end{proof}

If \(Z_o=0\), the proposition does not apply. In that case the metric
regularizer is Hessian-flat at the antipodal Dirac target by
\cref{eq:mf-center-EXP-regularizer}.
